\section{Discrete equivalent layers: The classical formulation}
\label{eq:equivalent_layer_technique}

Consider a set of 3D gravity or magnetic sources located in the subsurface,
over $\overline{\mathcal{V}}$ (see the subsection \textit{Potential fields due to 3D sources}). 
Let $\mathbf{d}^{o}$ be an $N \times 1$ observed data vector whose $i$-th element
$d^{o}_{i} \equiv d^{o}(\mathbf{r}_{i})$ (equation \ref{eq:observed-data-space}) 
represents a potential field produced by the 3D sources 
at an observation point $\mathbf{r}_{i} = ( x_{i} , y_{i} , z_{i} )^{\top} \notin \overline{\mathcal{V}}$
located above them, with $i \in \{ 1, \dots, N \}$.

Now, consider that $d^{o}_{i}$ can be approximated within a given accuracy $\epsilon$ 
by the potential field produced at $\mathbf{r}_{i}$ by a discrete and finite set of $M$ 
virtual sources such as monopoles (or point masses) or dipoles, for example, with 
constant and unknown physical property $p_{j}$ (e.g., mass, magnetic-moment intensity). 
These virtual sources are usually called \textit{equivalent sources}.
Here, we limit ourselves to the particular case in which they are disposed on a 
generic horizontal plane located below the potential-field data, at points 
$\mathbf{r}''_{j} = ( x''_{j} , y''_{j} , z''_{0} )^{\top}$,
$j \in \{ 1, \dots, M \}$, with constant vertical coordinate $z''_{0} > z_{i}$, $i \in \{ 1, \dots, N \}$.

Then, we may call the $M \times 1$ vector $\mathbf{p}_{\epsilon}$ with $j$-th element $p_{j}$ 
a \textit{discrete equivalent layer} if its predicted potential field 
\begin{equation}
	\mathbf{d}_{\epsilon} = \mathbf{A} \, \mathbf{p}_{\epsilon}
	\label{eq:predicted-data-vector}
\end{equation}
satisfies 
\begin{equation}
	\| \mathbf{d}^{o} - \mathbf{d}_{\epsilon} \|_{2}^{2} \le \epsilon \: ,
	\label{eq:epsilon-bound-data-discrete}
\end{equation}
where $\mathbf{d}_{\epsilon}$ is the predicted data vector whose $i$-th element represents 
the potential field due to $\mathbf{p}_{\epsilon}$ at the observation point $\mathbf{r}_{i}$ 
and $\mathbf{A}$ is an $N \times M$ matrix whose element $ij$ is defined by a harmonic function
$a_{ij} \equiv D_{a} \Psi(\mathbf{r}_{i}, \mathbf{r}''_{j})$, which represents a 
partial derivative of $\Psi$ (equation \ref{eq:Psi}), with $D_{a}$ given by the space-domain differential operator 
$D(\ell, \hat{\mathbf{u}})$ (equation \ref{eq:diff-operators-generic}).
The only conditions generally imposed on $a_{ij}$ are that they must 
(i) depend only on the coordinates of the $i$-th observation point and the $j$-th equivalent source and 
(ii) vanish as the distance between the $i$-th observation point and the $j$-th equivalent source increases.

After estimating a particular parameter vector $\tilde{\mathbf{p}}_{\epsilon}$ satisfying 
\ref{eq:epsilon-bound-data-discrete}, a potential field transformation such as upward continuation, 
interpolation or reduction to the pole, for example, can be computed as follows:
\begin{equation}
	\mathbf{f}_{\epsilon} = \mathbf{T} \, \tilde{\mathbf{p}}_{\epsilon} \: ,
	\label{eq:transformed-field-layer-space} 
\end{equation}
where $\mathbf{f}_{\epsilon}$ is an $L \times 1$ vector whose $k$-th element $f_{k}$ approximates 
the transformed potential field $f(\mathbf{r}^{\ddagger}_{k})$ (equation \ref{eq:transformed-field-space}) 
at a point $\mathbf{r}^{\ddagger}_{k} = (x^{\ddagger}_{k}, y^{\ddagger}_{k}, z^{\ddagger}_{k})^{\top}$, 
with $z^{\ddagger}_{k} < z''_{0}$, and $\mathbf{T}$ is an 
$L \times M$ matrix whose element $t_{kj}$ defines the transformed
potential field due to the $j$-th equivalent source, with $k \in \{ 1, \dots, L \}$ and
$j \in \{ 1, \dots, M \}$. 
We have defined the transformed potential field $f_{k}$ at generic points 
$\mathbf{r}^{\ddagger}_{k} \notin \overline{\mathcal{V}}$ above the layer plane $z''_{0}$, but they 
can coincide with those at which the observed potential field is measured. 

From the physical point of view, the product of the $j$-th element $p_{j}$ of $\mathbf{p}_{\epsilon}$ 
and the harmonic function $a_{ij}$ representing the element $ij$ of matrix $\mathbf{A}$ 
(equation \ref{eq:predicted-data-vector}) defines the potential field produced at the $i$-th observation
point by the $j$-th equivalent source with scalar physical property $p_{j}$.
For example, if $\mathbf{d}^{o}$ contains gravity disturbance data
\cite[e.g.,][ p. 93]{hofmann-wellenhof-moritz2005}, then
$p_{j} \, a_{ij}$ commonly defines the vertical component of the gravitational attraction
produced at $\mathbf{r}_{i}$ by a monopole located at
$\mathbf{r}''_{j}$, with mass $p_{j}$. 
This is the pioneer approach presented by \cite{dampney1969} for 
processing gravity data.
For the case in which $\mathbf{d}^{o}$ contains total-field anomaly data 
\cite[e.g.,][ p. 178]{blakely1996} instead, the common approach consists in
following \cite{emilia1973} and defining the term $p_{j} \, a_{ij}$ such
that it represents a magnetic induction field component 
produced at the $i$-th observation point by a dipole located at $\mathbf{r}''_{j}$,
with magnetic moment intensity $p_{j}$.

Note that, in the examples presented above, the physical relationship 
between the observed potential field data and their true sources is the
same as that between the data and the equivalent sources.
However, many authors 
\cite[e.g.,][]{cordell1992, guspi_novara2009, li_etal2014, soler_uieda2021} 
have proposed equivalent-layer methods by presuming that
the harmonic functions $a_{ij}$ are 
independent of the nature of the observed potential-field data and can be 
conveniently chosen according to the characteristics of the problem.

PAREI AQUI
